\section{System Model and Threat Assumptions}
\label{sec:03}
\label{sec:system-model}

\subsection{System Model and Participating Entities}
\label{subsec:system-entities}

The present study extends the biometric voting architecture reported in
\emph{Blockchain--integrated IoT biometric architecture for secure and
decentralized electronic voting} \citep{rifat2026biometric}. That study
combined fingerprint and facial verification with off-chain identity
management and blockchain vote processing. In the present system, the
authorization step does not receive the selected candidate. The voter-facing
client constructs an encrypted ballot, and the contract verifies its
zero-knowledge proof before accumulating ciphertexts. The application
maintains a Semaphore-based registration group and election anchor, whereas
the private-ballot pathway uses a separate custom Poseidon commitment tree.
The \texttt{VoteChainDirect} contract checks membership against the custom
eligibility root. These structures serve different purposes and do not form
a single membership proof. Biometric recognition results from the earlier
publication are not measurements of the present cryptographic procedure.

The participating entities are the voter-facing client and privacy worker,
biometric authentication service, credential authorization service,
election administrator, voting contract, public blockchain observers, five
tally trustees, and an external transcript auditor. The authentication
service identifies voters, while the credential service controls access
to registered voting secrets and releases authorization material. The
client prepares encrypted ballots and proofs. The contract validates
submitted ballots and maintains the encrypted aggregate. The administrator
configures election state and submits recovered results. Trustees provide
partial-decryption evidence, which an external auditor can independently
check. Public observers can inspect transactions, nullifiers, proofs, and
ciphertexts, but not plaintext candidate choices from accepted ballot data.
Administrative authority alone does not establish access to credential
secrets or trustee shares. These roles describe access boundaries without
assuming independent organizational control.

\subsection{Threat Model and Security Assumptions}
\label{subsec:system-security-objectives}

The threat model considers observers of public transaction data and
submitters attempting ineligible, malformed, or repeated ballots. It
also includes compromised voting clients, credential custodians, and
colluding trustees. Ballot acceptance depends on the configured circuit
relation and contract checks. A Groth16 proof binds membership in the
custom eligibility tree, the election context, a one-hot candidate
selection, and the submitted encrypted ballot slots
\citep{groth2016snark}. The contract verifies the proof and rejects a
previously used election nullifier before adding a ballot to the
encrypted tally. These checks do not independently establish correct
civil-identity enrollment, a fresh biometric match at transaction
submission, or the integrity of the voter's device.

Let $s$ be a registered private voting-credential secret and $e$ the
election identifier. The circuit computes the election-scoped nullifier
as
\begin{equation}
  n_e = \mathsf{Poseidon}(d_{\mathrm{null}},s,e).
  \label{eq:system-nullifier}
\end{equation}
where $d_{\mathrm{null}}$ denotes the configured domain separator.
Poseidon is used as a hash primitive within the proof circuit
\citep{grassi2021poseidon}. The nullifier is public, while the
membership witness and candidate selection are private. Nevertheless,
the credential authorization service can recover a stored credential
secret during authorization. It can therefore compute the corresponding
nullifier and link an identified voter to an accepted public transaction.
This identity linkage does not by itself reveal the candidate choice.
Candidate confidentiality also depends on ballot encryption and the
integrity of the client. A compromised client may expose the selection
before encryption.

Ballot confidentiality further depends on the protection of the
election decryption key shares. Three colluding trustees with sufficient
shares could decrypt individual published ballots, not only the
aggregate tally. The contract verifies ballot proofs on-chain, but
distributed-key-generation and partial-decryption evidence are
verified externally. Finalization records reported totals and a
transcript hash after contract-state and arithmetic checks. It does not
verify trustee decryption proofs on-chain. Independent tally verification
requires access to the full transcript and its corresponding encrypted
aggregate. These assumptions do not establish receipt-freeness,
coercion resistance, or security under production-election conditions.
